\chapter{结果展示示例}

\section{结果描述示例}

在实验或实现章节中，通常需要配合图表对结果进行定量描述，并对现象作出尽量客观的解释。对于需要展示多项指标的结果，也可以通过三线表进行归纳。

\begin{bupttable}{示例结果表}{tab:sample_results}
    \begin{tabular}{@{}p{2.8cm} p{2.8cm} p{2.8cm} p{3.6cm}@{}}
        \toprule
        条件 & 指标一 & 指标二 & 说明 \\
        \midrule
        配置 A & 0.91 & 1.25 & 仅用于模板排版示例 \\
        配置 B & 0.85 & 1.10 & 仅用于模板排版示例 \\
        配置 C & 0.72 & 0.96 & 仅用于模板排版示例 \\
        \bottomrule
    \end{tabular}
\end{bupttable}

\section{讨论示例}

结果讨论部分应围绕数据展开，避免无依据的夸张性描述。若某项指标优于其他配置，应说明其可能原因；若结果不如预期，也应指出潜在限制和后续改进方向。

\section{本章小结}

本章说明了结果展示与讨论的基本写法，可据此扩展为正式论文中的实验分析章节。
